% paper/sections/10c_fccd_nonuniform.tex
%
% §10.4  非一様格子上の FCCD
%         §4.6 で定義した一様格子上の FCCD を非一様格子へ拡張する．
%
% ═══════════════════════════════════════════════════════════════════════════════
\subsection{非一様格子上の FCCD：\texorpdfstring{$\mathcal{O}(H^3)$}{O(H3)} 行列定式（\texorpdfstring{$\theta=1/2$}{θ=1/2} 界面適合で \texorpdfstring{$\Ord{H^4}$}{O(H4)}）}
\label{sec:fccd_nonuniform}
% ═══════════════════════════════════════════════════════════════════════════════

\S\ref{sec:fccd_matrix} で導入した一様格子版行列
$\mathbf{M}^{\FCCD} = \mathbf{D}_1 - \mathbf{D}_2\mathbf{S}_{\CCD}$ を
非一様格子に拡張するに際し（コンパクト系の非一様格子拡張については~\cite{Lele1992}を再掲），
非一様 \textbf{打ち消し係数} $(\mu_i,\lambda_i)$ を
3 本の疎な双対角行列に織り込む．

% -------------------------------------------------------------------------------
\subsubsection{幾何量と打ち消し係数}
\label{sec:fccd_nu_geometry}
% -------------------------------------------------------------------------------

\S\ref{sec:grid_density} の格子密度関数が定めた格子において，
節点を $x_0 < x_1 < \cdots < x_N$ とし，
内部面を $f_{i-1/2}$（$i=1,\ldots,N$）と記す．
局所幾何量を：
\begin{align}
  H_i &:= x_i - x_{i-1} \quad\text{（面幅）}, \label{eq:fccd_nu_Hi}\\
  \theta_i &:= (x_i - x_{f_{i-1/2}})/H_i \quad\text{（左節点値の線形補間重み）}. \label{eq:fccd_nu_theta}
\end{align}
中央（セル中心面）では $\theta_i \equiv 1/2$．
\textbf{打ち消し係数}：
\begin{equation}
  \mu_i := \theta_i - \tfrac{1}{2},
  \qquad
  \lambda_i := \frac{1 - 3\theta_i(1-\theta_i)}{6}.
  \label{eq:fccd_nu_mulambda}
\end{equation}
$\mu_i$ は面の中央ズレ，$\lambda_i$ は $u_f'''$ 近似の非一様補正項を担う．
一様極限 $\theta_i\to 1/2$ で $\mu_i\to 0$，$\lambda_i\to 1/24$．

% -------------------------------------------------------------------------------
\subsubsection{3 本の疎な双対角行列}
\label{sec:fccd_nu_matrices}
% -------------------------------------------------------------------------------

$\mathbb{R}^{N\times(N+1)}$ 上の
3 本の疎な双対角行列 $\mathbf{D}_1^{(H)}, \mathbf{D}_\mu^{(H\theta)}, \mathbf{D}_\lambda^{(H)}$ を定義する．
行 $f_{i-1/2}$ は 2 つの節点列（$i-1$ と $i$）のみ非零：
\begin{align}
  (\mathbf{D}_1^{(H)})_{f_{i-1/2},i-1} &= -1/H_i,
  &\quad (\mathbf{D}_1^{(H)})_{f_{i-1/2},i} &= +1/H_i, \label{eq:fccd_nu_D1}\\
  (\mathbf{D}_\mu^{(H\theta)})_{f_{i-1/2},i-1} &= \mu_i H_i\theta_i,
  &\quad (\mathbf{D}_\mu^{(H\theta)})_{f_{i-1/2},i} &= \mu_i H_i(1-\theta_i), \label{eq:fccd_nu_Dmu}\\
  (\mathbf{D}_\lambda^{(H)})_{f_{i-1/2},i-1} &= -\lambda_i H_i,
  &\quad (\mathbf{D}_\lambda^{(H)})_{f_{i-1/2},i} &= +\lambda_i H_i. \label{eq:fccd_nu_Dlambda}
\end{align}
$\mathbf{D}_1^{(H)}$ は 1 次差分，$\mathbf{D}_\mu^{(H\theta)}$ は $u_f''$ 補間，
$\mathbf{D}_\lambda^{(H)}$ は $u_f'''$ 差分に対応する．

% -------------------------------------------------------------------------------
\subsubsection{非一様 FCCD 行列}
\label{sec:fccd_nu_matrix_form}
% -------------------------------------------------------------------------------

非一様 FCCD の面式
\begin{equation}
  d_{f_{i-1/2}} = \frac{u_i - u_{i-1}}{H_i}
    - \mu_i H_i\bigl[\theta_i q_{i-1} + (1-\theta_i)q_i\bigr]
    - \lambda_i H_i(q_i - q_{i-1})
  \label{eq:fccd_nu_face_formula}
\end{equation}
をベクトル形で読み替えると：
\begin{equation}
  \mathbf{d}^{\FCCD,\text{nu}}
  = \mathbf{D}_1^{(H)}\mathbf{u}
   - (\mathbf{D}_\mu^{(H\theta)} + \mathbf{D}_\lambda^{(H)})\mathbf{q},
\end{equation}
したがって節点から面への非一様 FCCD 演算子は
\begin{equation}
  \boxed{\;
  \mathbf{M}^{\FCCD,\text{nu}}
  = \mathbf{D}_1^{(H)}
    - \bigl(\mathbf{D}_\mu^{(H\theta)} + \mathbf{D}_\lambda^{(H)}\bigr)
      \mathbf{S}_{\CCD}.
  \;}
  \label{eq:fccd_nu_composite_matrix}
\end{equation}
ここで $\mathbf{S}_{\CCD}$ は \S\ref{sec:fccd_matrix} で定義した
非一様 CCD ブロック三重対角系の第二導関数射影であり，
\S\ref{sec:ccd_metric} の座標変換により非一様格子上で直接構築できる．

\noindent\textbf{一様極限（検証）}：
$\theta_i \equiv 1/2 \Rightarrow \mu_i = 0$，$\lambda_i = 1/24$，
$\mathbf{D}_\mu^{(H\theta)} \equiv \mathbf{0}$，
$\mathbf{D}_\lambda^{(H)} = \mathbf{D}_2$（\S\ref{sec:fccd_matrix} 式~\eqref{eq:fccd_D2}）．
したがって式~\eqref{eq:fccd_nu_composite_matrix} は
\S\ref{sec:fccd_matrix} 式~\eqref{eq:fccd_composite_matrix} と一致する．

% -------------------------------------------------------------------------------
\subsubsection{切断誤差：\texorpdfstring{$\mathcal{O}(H^3)$}{O(H3)}}
\label{sec:fccd_nu_truncation}
% -------------------------------------------------------------------------------

Taylor 解析（式~\eqref{eq:fccd_nu_mulambda} の $\mu,\lambda$ 導出と整合）により，
一般の $\theta_i \neq 1/2$ に対する FCCD 切断誤差は
\begin{equation}
  D^{\FCCD,\text{nu}} u_{f_{i-1/2}}
  = u'_{f_{i-1/2}} + c_3(\theta_i) H_i^3 u^{(4)}_{f_{i-1/2}} + \Ord{H_i^4}
  \label{eq:fccd_nu_truncation}
\end{equation}
となり，セル中心面以外では $\Ord{H^3}$ に落ちる．
$c_3(\theta) = (\theta-1/2)(1-\theta(1-\theta))/24$ は
$\theta = 1/2$ で厳密に $0$ となるので，
セル中心格子では $\Ord{H^4}$ が維持される．
本稿の二相流標準構成では面位置を常に中心（$\theta = 1/2$）にとる
\textbf{界面適合格子族}（\S\ref{sec:grid_density}）を採用するため，
面演算子としての局所次数は $\Ord{H^4}$ である．
ただし物理座標 $x$ への写像を台形則で構成する場合，
合成された物理空間微分の全体次数は座標生成誤差に律速される
（§\ref{sec:grid_metric_order_note} 参照）．
壁面近傍，境界補助面，HFE 最近接点，GFM 交差点，ジャンプ補正面で
$|\theta_i-1/2|>\theta_{\mathrm{tol}}$ となる場合は，
その面を $\Ord{H^3}$（または境界閉包に応じて $\Ord{H^2}$）の特殊面として扱い，
$\Ord{H^4}$ の検証集合から分離する．

% -------------------------------------------------------------------------------
\subsubsection{幾何重みと演算子構造}
\label{sec:fccd_nu_geometry_cost}
% -------------------------------------------------------------------------------

3 本の重み行列 $\mathbf{D}_1^{(H)}, \mathbf{D}_\mu^{(H\theta)}, \mathbf{D}_\lambda^{(H)}$ は
格子幾何 $(H_i,\theta_i)$ のみに依存するため，
固定格子では時間に依存しない幾何係数として定まる．
各行は隣接する 2 節点のみに結合し，非一様化は係数行列の値を変えるだけで
FCCD 面演算子の局所結合構造を変えない．
したがって，一様版と非一様版は同じ節点--面対応の上で比較できる．

% -------------------------------------------------------------------------------
\subsubsection{Balanced--Force への影響}
\label{sec:fccd_nu_bf_impact}
% -------------------------------------------------------------------------------

\S\ref{sec:fccd_face_residual_order} で示した BF 残差 $|\BF_\text{res}| = \Ord{\Delta x^4}$ の性質は，
非一様格子でも界面適合（$\theta=1/2$）を維持する限り保たれる．
界面近傍のみ $\theta=1/2$ を保証する「局所中央化」を
\S\ref{sec:grid_density} の格子密度関数で課すことにより，
BF 原理（\S\ref{sec:bf_operator_mismatch}）の要求する
$p$ 勾配と CSF の同一面位置は保たれる．
ただしこれは滑らかな場を同じ FCCD 面演算子に載せる場合の主張であり，
圧力ジャンプ/HFE 型のジャンプ補正閉包では，§~\ref{sec:split_ppe_pressure_jump_formulation} の
$B_f^{\mathrm{nu}}=s_fj_f/H_f$ と同一界面面データを PPE 右辺と速度補正で共有して初めて
BF 条件が閉じる．
非一様メトリクス $G_f^{\mathrm{nu}}$ だけでは，
既知ジャンプ $B_f^{\mathrm{nu}}$ は自動的には生成されない．

\noindent 詳細な離散化補足・PCR/Thomas 法の非一様形・
境界行の壁面境界条件拡張は付録~\ref{app:ccd_boundary_alt} 参照．

% ═══════════════════════════════════════════════════════════════════════════════
